\section{Sectorial realizations, material composition, and collection spectra}
\label{vi:sec:familias-compuestos}

The atlas of Section~\ref{vi:sec:atlas} and the reader of
Section~\ref{vi:sec:operador-masa} have complementary roles. The former
classifies supports, depths, and orbits; the latter preserves the chronology
of each trajectory and expresses its mass operator. Applying them to a collection
requires both structures to be retained. Equality between mass values does not
identify routes; equality between cells does not identify histories,
representations, or composite states either.

This section develops the sectorial realizations of that construction.
The internal constants, angular pair, absolute scale, and bidirectional readers
are those already constructed from APP--TRIT--TPK. The physical representation
specifies which information of the state is observed; metrological comparison
takes place after that representation has been fixed. In particular, the
exposition preserves the distinction between a fiber operator, an exact evaluation
for a declared signature, and a spectral line individuated by a coupling.

\subsection{Collection descriptors and the spectral criterion for individuation}
\label{vi:subsec:descriptor-familia}

For an elementary state, write
\[
 X=(\mathfrak c,\rho,g,f,q,\epsilon,\varsigma,\mathfrak a).
\]
Here \(\mathfrak c\) is the sectorial class, \(\rho\) the representation,
\(g\) the physical generation, \(f\) flavor, \(q\) charge,
\(\epsilon\) sheet orientation, \(\varsigma\) statistical parity,
and \(\mathfrak a\) the coupling data. The combinatorial depth
of the atlas and the physical generation remain coordinates of different
types. Flavor distinguishes components of a collection representation;
color distinguishes components of the color representation;
the sheet and its orientation preserve the operatorial provenance of the
trajectory. None of these labels replaces the event register.

Let \(\mathcal P_{\rm adm}\) be the set of admissible physical descriptors,
let \(\mathcal R_{\rm adm}\) be the register of enriched routes produced by
the atlas, and let \(\mathcal D\) be the space of incidence descriptors that
preserve sectorial class, representation, depth, charge, color,
orientation, and coupling data. Define
\[
 d_{\mathcal P}:\mathcal P_{\rm adm}\longrightarrow\mathcal D,
 \qquad
 d_{\mathcal R}:\mathcal R_{\rm adm}\longrightarrow\mathcal D,
\]
where \(d_{\mathcal P}(X)\) expresses the structural requirements of state
\(X\), and \(d_{\mathcal R}(\gamma)\) expresses the same coordinates preserved
by route \(\gamma\). The fiber, orbit, or multisection produced by the
descriptor is the fiber product
\begin{equation}
 F_X=d_{\mathcal R}^{-1}\!\bigl(d_{\mathcal P}(X)\bigr)
 =\mathcal R_{\rm adm}\times_{\mathcal D}\{d_{\mathcal P}(X)\}.
 \label{vi:eq:fibra-descriptor}
\end{equation}
For the realizable descriptors in this article, \(F_X\) is finite and nonempty.
On
\(\mathcal H_X=\ell^2(F_X)\otimes V_{\rho_X}\) act the two operators
\(M_X^{\beta,D}\) and \(M_X^{\beta,\Omega}\) of
Section~\ref{vi:sec:operador-masa}. Since they are diagonal in the route
basis, their joint spectral measure is
\begin{equation}
 E_X(B)=\sum_{\gamma\in F_X}
 \mathbf1_B\bigl(m_D(\gamma),m_\Omega(\gamma)\bigr)
 |\gamma\rangle\langle\gamma|\otimes I_{V_{\rho_X}},
 \qquad B\subset\mathbb R_{>0}^2.
 \label{vi:eq:familias-pvm}
\end{equation}
Preserving this measure allows two fibers to be distinguished even when their means
coincide. A scalar functional of the two readers is applied only
when it forms part of the declared representation. If the coupling
preserves several values, the output is a spectrum with multiplicities.

\begin{proposition}[Spectral-line criterion]
\label{vi:prop:familias-linea}
Let \(M\) be a self-adjoint operator on a finite fiber and \(P\) a
preparation projector. The channel expresses exclusively the value \(m\)
if and only if
\[
 E_M(\{m\})P=P,
 \qquad\text{equivalently}\qquad (M-mI)P=0.
\]
The equality \(PMP=mP\) alone determines a scalar compression,
but does not guarantee that the prepared spectral support is a single point.
\end{proposition}
\begin{proof}
By the spectral theorem, \(\ker(M-mI)\) is the image of
\(E_M(\{m\})\). The first two conditions state that
\(\operatorname{im}P\) is contained in that eigenspace. To verify
the last distinction, take \(M=\operatorname{diag}(1,3)\) and the projector
onto \((1,1)/\sqrt2\). Then \(PMP=2P\), whereas the prepared
mass distribution has support \(\{1,3\}\), not \(\{2\}\).
\end{proof}

Individuation thus preserves the scope of each result. The fiber
and operator are determined objects; passage to a line uses the
projector and spectral condition stated above. This formulation will apply
equally to charged, neutral, and composite collections.

\subsection{Charged leptons: electronic center and generation multisections}
\label{vi:subsec:leptones-cargados}

The three charged multisections have ranks \(1,8,16\) at depths
\(G_0,G_2,G_4\). They share the charge representation and trivial
color, but not the support rank or trajectory register.
The electron occupies the center \((5,5)\); the muonic and
tauonic realizations preserve their multisections until spectral
coupling is performed. Conjugation of their representations is treated through the
operators of Section~\ref{vi:sec:conjugacion}, not through a new
choice of masses.

The basal electronic coordinate \(\mathfrak e_0\) provides the intermediate
angular chart
\begin{equation}
 \mathfrak m_X^{(0)}=\mathfrak e_0
 \exp\left(n_A(X)x+\frac{s_C(X)}6y\right),
 \qquad x=A_{\mathrm{rad}},\quad y=C^*_{\mathrm{rad}}.
 \label{vi:eq:carta-angular-familias}
\end{equation}
In this formula, \(\mathfrak e_0\) is the basal section, not the beta
electronic closure. Bidirectional return belongs to the subsequent stage
of the construction. The two stages preserve their own normalization.

The recovered nominal registers for the two noncentral collections
express
\[
 (n_A,s_C)_\mu=(42,-3),\qquad (n_A,s_C)_\tau=(64,0),
\]
and their declared complete signatures are
\begin{equation}
 \sigma_\mu=(42,-3,8,0,2),\qquad
 \sigma_\tau=(64,0,25,-1,6).
 \label{vi:eq:firmas-muon-tau}
\end{equation}
To explain the effect of each coordinate, denote by \(\mathcal E_j\)
the successive readings of a declared signature:
\begin{align}
 \mathcal E_0(\sigma)&=n_Ax+s_Cy/6,\notag\\
 \mathcal E_1(\sigma)&=\mathcal E_0(\sigma)+k_{\Delta_4}\Delta_4,\notag\\
 \mathcal E_2(\sigma)&=\mathcal E_1(\sigma)+b_{120}s_{120},\notag\\
 \mathcal E_3(\sigma)&=\mathcal E_2(\sigma)+b_\zeta\zeta_{270},
 \qquad\zeta_{270}=135\Delta_4^2.
 \label{vi:eq:etapas-e3-sectoriales}
\end{align}
The symbol \(\mathcal E_3\) here denotes an evaluation exponent,
not a species or an additional mode. The beta supplement preserving these
signatures uses the explicit normalization
\(m_X^{[3]}=m_e^\beta\exp\mathcal E_3(\sigma_X)\).
This is an evaluation relative to the supplement's signatures; it does not replace
the route-by-route construction of \(M_X^{\beta,D}\) and
\(M_X^{\beta,\Omega}\).

The difference of signatures gives, without approximation,
\begin{equation}
 \log\frac{m_\tau^{[3]}}{m_\mu^{[3]}}
 =22x+\frac12y+17\Delta_4-s_{120}+4\zeta_{270}.
 \label{vi:eq:razon-tau-muon}
\end{equation}
The proof consists of subtracting the exponents in
\eqref{vi:eq:etapas-e3-sectoriales}. The angular term, torsional
memory, and two return readings remain identified in the
result. The preserved evaluations are
\begin{align*}
 m_\mu^{[3]}&=105.658360294435829303\ldots\ \mathrm{MeV}/c^2,\\
 m_\tau^{[3]}&=1776.860917631432295595\ldots\ \mathrm{MeV}/c^2.
\end{align*}
Their documentary status is exact evaluation upon fixing the recovered signature.
Physical selection of that signature within the multisection and the spectral-line
criterion are distinct operations and are preserved as such.

At the electronic center the rank is one, and the readers coincide at
\((80,54,6)\). The active history
\((2,2,5,-4,-3,-2)^2\) preserves more information than its historical
envelope \((4,3,5,-4,-3,-5)^2\), although both share the sign
pattern. This distinction links the electronic particularity to the
positional memory of events: collections are not defined solely by
the visual form of a sign word.

\subsection{Neutrinos: neutrality, mass spectrum, and flavor transport}
\label{vi:subsec:neutrinos-familias}

The neutrino support has dimension \(20=2+4+6+8\), distributed over
four depths. Its three flavor labels
\(\nu_e,\nu_\mu,\nu_\tau\) and its three possible mass eigenvalues
belong to the physical realization of the sector, not to the census of
depths. Electrical neutrality preserves orientation,
sheet memory, and spinorial holonomy. It therefore does not eliminate the
variables involved in mixing.

The neutral register separates the contributions of two half-cycles. If
\(\mathcal I_m\) is the block of \(9\) iterations of phase \(m\),
\(\theta_t\) marks the visit to \(9\), \(\epsilon_t\) preserves the sheet,
and \(\omega_t\) its corona orientation, the local reader is
\[
 T_\nu(m)=\sum_{t\in\mathcal I_m}\omega_t\epsilon_t\theta_t,
 \qquad e_m^\nu=T_\nu(m)-T_\nu(m+6).
\]
\subsubsection{Exact cancellations of the antipodal extractor}
\label{vi:subsubsec:cancelaciones-antipodales}

Indices in this subsection are read in \(\mathbb Z/12\mathbb Z\), and
\(\operatorname{sgn}(0)=0\). Let \(A_r\) be the phase shift
\((A_rT)(m)=T(m+r)\). For the neutral register,
\begin{equation}
 e^\nu=(I-A_6)T_\nu,
 \qquad e^\nu_{m+6}=-e^\nu_m.
 \label{vi:eq:extractor-antipodal}
\end{equation}
Define
\[
 B_a^\nu=e_a^\nu+e_{a+4}^\nu+e_{a+8}^\nu,
 \qquad
 \tau_m^\nu=\operatorname{sgn}(e_m^\nu),
 \qquad
 c_m^\nu=\tau_m^\nu\tau_{m+3}^\nu.
\]

\begin{proposition}[Vanishing of the two antipodal scalar readers]
\label{vi:prop:cancelaciones-antipodales}
The neutral extractor satisfies
\begin{equation}
 B_{a+2}^\nu=-B_a^\nu,\qquad
 b_{120}^\nu:=\sum_{a=0}^{3}\operatorname{sgn}(B_a^\nu)=0,
 \label{vi:eq:cancelacion-b120}
\end{equation}
and
\begin{equation}
 c_{m+3}^\nu=-c_m^\nu,\qquad
 b_\zeta^\nu:=\sum_{m=0}^{11}c_m^\nu=0.
 \label{vi:eq:cancelacion-bzeta}
\end{equation}
\end{proposition}

\begin{proof}
By~\eqref{vi:eq:extractor-antipodal},
\[
 B_{a+2}^\nu
 =e_{a+2}^\nu+e_{a+6}^\nu+e_{a+10}^\nu
 =-e_{a+8}^\nu-e_a^\nu-e_{a+4}^\nu=-B_a^\nu.
\]
The four summands in \(b_{120}^\nu\) cancel in the pairs
\((0,2)\) and \((1,3)\), including when one of them vanishes. Likewise,
\(\tau_{m+6}^\nu=-\tau_m^\nu\), whence
\[
 c_{m+3}^\nu
 =\tau_{m+3}^\nu\tau_{m+6}^\nu
 =-\tau_{m+3}^\nu\tau_m^\nu=-c_m^\nu.
\]
The sum of the twelve terms vanishes in alternating pairs.
\end{proof}

Scalar cancellation does not destroy oriented information. On the
anti-invariant subspace
\[
 S_-:=\ker(A_6+I),
 \qquad J:=A_3\big|_{S_-},
\]
one has \(J^2=-I\). Since the shift is unitary,
\(J^*=J^{-1}=-J\): the register retains an oriented complex structure.
Moreover,
\[
 e^\nu=(I-A_6)T_\nu=2P_-T_\nu,
 \qquad P_-=(I-A_6)/2,
\]
so the complete antipodal vector remains available even though its two
scalar contractions vanish.

The cyclic return \(A_{12}=I\) concerns this phase representation; it is
not the return of the complete TPK state, which retains carry, sheet,
boundary, and memory. For the same reason,
\eqref{vi:eq:cancelacion-b120}--\eqref{vi:eq:cancelacion-bzeta} do not
annul the general coefficients \(b_{120}\) and \(b_\zeta\) of a
non-antipodal route, nor do they modify masses, composite signatures, or
PMNS transport. The longitudinal seam used in register prolongation is
controlled independently by Lemma~\ref{vi:lem:costura-108-1080}.


This difference retains the oriented information even when its total
sum vanishes. The complete operator also uses torsional
predecessors, vacancies, and transport between sheets. Passage to the
physical subspace requires the rank-\(3\) neutral projector,
\[
 M_\nu=P_\nu\mathcal M_\nu P_\nu
 \quad\text{on }\operatorname{im}P_\nu.
\]
The eigenvalues of that compression are the masses in the internal
action--clock--mass scale. Mixing angles act between their frames;
they neither replace the scale nor retrospectively generate the eigenvalues.
The compression coincides with the restriction of the complete operator when
\(\operatorname{im}P_\nu\) is a reducing subspace of
\(\mathcal M_\nu\).

The register allows a Gram kernel to be assembled before diagonalization.
For each data type
\(r\in\{\mathrm{car},\mathrm{hoj},\mathrm{ori},\mathrm{vac},\mathrm{hol}\}\),
let \(\Xi_r(c)\) be the vector produced by cell or route \(c\).
Define
\begin{equation}
 K_{\mathrm{reg}}(c,d)=
 \sum_{r:N_r>0}\frac{\langle\Xi_r(c),\Xi_r(d)\rangle}{N_r},
 \qquad N_r=\sum_a\|\Xi_r(a)\|^2.
 \label{vi:eq:gram-neutrinos}
\end{equation}
Omission of a term with zero total norm is fixed before
evaluation. It preserves a well-defined domain and introduces no fit.

\begin{proposition}[Positivity and rank of the neutral register]
\label{vi:prop:gram-neutrinos}
The kernel \(K_{\mathrm{reg}}\) is Hermitian and positive semidefinite.
For a three-dimensional aggregation \(G\), the condition \(\det G\ne0\)
certifies that the aggregated data separate three directions. For any
invertible square matrix \(G\), its polar factor
\[
 U_G=G(G^*G)^{-1/2}
\]
is unitary. This factor does not by itself determine the spectral frames
of the charged and neutral operators.
\end{proposition}
\begin{proof}
For coefficients \(z_c\), the quadratic form of
\eqref{vi:eq:gram-neutrinos} is
\[
 \sum_{c,d}\overline z_cK_{\mathrm{reg}}(c,d)z_d
 =\sum_{r:N_r>0}N_r^{-1}
 \left\|\sum_c z_c\Xi_r(c)\right\|^2\ge0.
\]
The rank is the dimension of the image of the aggregated vectors. For the
last statement,
\(U_G^*U_G=(G^*G)^{-1/2}G^*G(G^*G)^{-1/2}=I\);
invertibility and equality of dimensions also give \(U_GU_G^*=I\).
The formula uses matrix \(G\), whereas a spectral frame requires
diagonalization of the corresponding mass operator.
\end{proof}

If \(F_\ell,F_\nu:\mathbb C^3\to\mathcal H\) are complete orthonormal
frames of the same physical subspace, then
\(U_{\mathrm{PMNS}}=F_\ell^*F_\nu\) is unitary. If their images are
different, the product is only a contraction; unitarity requires
equality of the transported subspaces. In the presence of degeneracies,
mixing is determined as a double coset under frame changes in
the degenerate blocks, not as a numerical matrix chosen by convention.

The minimal phase construction
\[
 \Phi_K=I+(e^{i\pi/6}-1)|b_K\rangle\langle b_K|,
 \qquad\|b_K\|=1,
\]
remains a documented negative control: it is unitary, but the best
permutation of its angular identification gives
\(\theta_{13}\simeq17.02^\circ\), incompatible with the comparison for that
reduced model. The result concerns the rank-one phase. It does not
replace the complete kernel \eqref{vi:eq:gram-neutrinos}, nor permit a mixing
matrix that has not passed its own rank and angle checks
to be presented as validated.

\subsubsection{Neutral self-conjugation and the possibility of a sterile sector}

The Dirac--Majorana alternative uses the fermionic representation. A
neutral real structure must satisfy
\[
 \mathcal C_\nu^2=I,\qquad
 \mathcal C_\nu M_\nu=M_\nu\mathcal C_\nu,\qquad
 \mathcal C_\nu H_\nu=H_\nu\mathcal C_\nu,
\]
with \(\mathcal C_\nu\) antiunitary and a compatible realification.
Commutation with mass preserves its real part. For the self-adjoint
Hamiltonian, by contrast, antiunitarity also conjugates the factor
\(i\): if \(U_\nu(t)=\exp(-itH_\nu/\hbar)\), then
\[
 \mathcal C_\nu U_\nu(t)\mathcal C_\nu^{-1}=U_\nu(-t).
\]
Consequently, the fixed modes \(\psi=\mathcal C_\nu\psi\) define a
neutral real structure, but the stated conditions are insufficient to
restrict the entire Schr\"odinger evolution to it. Preservation of
that fixed part requires \(\mathcal C_\nu\) to commute with the group
\(U_\nu(t)\); Proposition~\ref{vi:prop:conjugacion-evolucion}
proves the corresponding condition on its generator. A
Majorana realization must also declare the compatibility of
self-conjugation with the fermionic representation and its dynamics.
A nonzero conserved charge that changes
sign under \(\mathcal C_\nu\) keeps conjugate modes separate
in the Dirac realization. The word involution of
Section~\ref{vi:sec:conjugacion} and electrical neutrality, considered
in isolation, do not decide this alternative.

A sterile sector requires a nonzero projector \(P_s\) with
\(P_sP_{\mathrm{act}}=0\) and \(P_sJ_{\mathrm{weak}}=0\).
If \([P_s,M_\nu]=0\), it is reducing and does not mix with the active sector.
If \(P_sM_\nu P_{\mathrm{act}}\ne0\), mixing exists and the projector
ceases to commute with mass. Consequently, an additional atlas depth
does not automatically count as a fourth neutrino species.

\subsection{Quarks: flavor, color orbit, and scale transport}
\label{vi:subsec:quarks-familias}

The quark state preserves the tuple
\[
 (\sigma_{ud},g,[\gamma]_{\mathrm{col}},\mathfrak M,\mathscr L),
\]
which distinguishes up or down type, generation, color orbit,
multisection, and register. Section~\ref{vi:sec:atlas} constructs the
color residue \(\kappa_{\mathrm{col}}(r,c)=r-c\pmod3\) and the
three classes of \(12\) cells. The color index acts within the
support; flavor preserves generation and signature information that this
residue does not identify.

The six nominal angular pairs are
\[
\begin{array}{c|r r|r r}
 f&n_A&s_C&W_+&W_-\\\hline
 u&11&7&73&59\\
 d&17&8&110&94\\
 c&61&8&374&358\\
 s&41&-3&243&249\\
 t&100&-1&599&601\\
 b&71&-5&421&431
\end{array}
\qquad W_\pm=6n_A\pm s_C.
\]
These last integers rewrite the same exponent:
\[
 n_Ax+s_Cy/6
 =\frac{W_+(x+y)+W_-(x-y)}{12}.
\]
Adding and subtracting \(W_\pm\) recovers \(n_A\) and \(s_C\), so that
this presentation preserves both channels. The six pairs belong
to the documented angular chart; the three remaining coordinates are not
filled with zeros to manufacture a complete signature. Their color
operators and angular evaluations are preserved as results of
different scope.

\begin{proposition}[Color degeneracy of a flavor mass]
\label{vi:prop:masa-color}
Let \(V_{\mathrm{col}}\simeq\mathbb C^3\) be the fundamental color
representation and \(U\) its irreducible unitary action. The reduction
\[
 \mathbb E_{\mathrm{col}}(M)=
 \int_{SU(3)}U(g)MU(g)^*\,dg
\]
is a multiple of \(I_3\). If \(M\) already commutes with the action, its three
color components have the same flavor-mass value.
\end{proposition}
\begin{proof}
Invariance of Haar measure gives
\(U(h)\mathbb E_{\mathrm{col}}(M)U(h)^*=
\mathbb E_{\mathrm{col}}(M)\). Schur's lemma makes this
operator scalar. Since trace is preserved, the scalar is
\(\operatorname{Tr}M/3\). If \(M\) is invariant from the outset,
the integral returns \(M\).
\end{proof}

The proof explains what the color quotient eliminates: the dependence of
the numerical value on the red, green, or blue representative. It does not eliminate the
color routes subsequently involved in a bound state. Nor does it
turn this reduction into a flavor selector.

The current quark mass is presented as a section at a declared scale and in a declared
scheme, \(m_q(\mu,\mathscr R)\). If \(t=\log\mu\) and the connection
induces \(\gamma_m(g(t))\), transport satisfies
\[
 \left(\frac{d}{dt}+\gamma_m(g(t))\right)m_q=0,
 \qquad
 m_q(t_2)=\mathcal P\exp\left(-\int_{t_1}^{t_2}
 \gamma_m(g(t))\,dt\right)m_q(t_1).
\]
The connection and its endomorphism are transport data, not numbers
chosen to match a table. A physical residual is meaningful only
after both sections have been placed at the same scale and in the same scheme.

CKM mixing acts on another coordinate: the flavor frame. Once
the internal unitary matrix has been constructed, the down operator, expressed
in the up frame, is \(B_d^{(u)}=V_{\mathrm{CKM}}B_dV_{\mathrm{CKM}}^*\).
Unitary conjugation preserves the characteristic polynomial, since
\[
 \det(zI-VB_dV^*)=\det\bigl(V(zI-B_d)V^*\bigr)=\det(zI-B_d).
\]
The role of the angles is thus understood: they relate frames without
becoming retrospective generators of mass eigenvalues.

\subsection{Zero-mass charts: the photon and the adjoint color representation}
\label{vi:subsec:nulo-familias}

The null domain is separated before the positive character is applied. In the
total realization,
\[
 \mathcal H_{\mathrm{phys}}=\mathcal H_0\oplus\mathcal H_+,
 \qquad M|_{\mathcal H_0}=0,\qquad M|_{\mathcal H_+}>0.
\]
The exponential of a finite signature does not vanish. Thus a null chart
is not obtained by assigning an arbitrarily large signature to a positive route
and calling its approximation zero.

The photon belongs to the electromagnetic chart, with its gauge
representation and its two physical transverse polarizations. Gluons
belong to the adjoint color representation, of dimension \(8\).
Equality of their zero masses does not identify the two representations.
The fiber, gauge, polarizations, and coupling channels
remain part of the state.

Oriented propagation with zero mass is also compatible. In the
documented chart,
\[
 c_\pm=\frac{c}{1\pm\varepsilon},\qquad
 \varepsilon=q_-^{90}-q_+^{90},\qquad0<\varepsilon<1.
\]
Then \(c_\pm>0\) and
\[
 \frac{2}{c_+^{-1}+c_-^{-1}}=c.
\]
The identity follows by adding \((1+\varepsilon)/c\) and
\((1-\varepsilon)/c\). The asymmetry of the two readings preserves its
orientation information; it does not by itself express a mass term.

\subsection{Massive bosons and the scalar sector}
\label{vi:subsec:bosones-familias}

The descriptors of \(W\), \(Z\), and the scalar sector are not forced
into one of the \(13\) fermionic multisections. They preserve their
representation and realization domain. The documented nominal
readings are
\[
\begin{array}{c|c|c|c|c}
 X&\text{nominal route}&\text{direction}&(n_A,s_C)&(W_+,W_-)\\\hline
 W&15/37&E\text{--}O&(94,-1)&(563,565)\\
 Z&20/23&E\text{--}O&(95,-1)&(569,571)\\
 H&24/29&N\text{--}S&(97,9)&(591,573)
\end{array}
\]
These route names belong to the preserved nominal register. Their
joining to the complete oriented domain must retain the trajectory,
representation, and projector data. The expressed cardinal counts,
respectively \((23,44,11,30)\), \((33,45,10,20)\), and
\((40,34,22,12)\), are preserved together with the direction and are not
replaced by a single label.

For \(W\) and \(Z\), the complete signatures are recovered:
\[
 \sigma_W=(94,-1,0,-2,4),\qquad
 \sigma_Z=(95,-1,-11,0,-4).
\]
Reading \eqref{vi:eq:etapas-e3-sectoriales} produces, for these signatures,
\[
 m_W^{[3]}=80378.964909704547024865\ldots\ \mathrm{MeV}/c^2,
\]
\[
 m_Z^{[3]}=91187.533444313407844855\ldots\ \mathrm{MeV}/c^2.
\]
Their ratio preserves exactly the difference of the evaluated registers:
\[
 \log\frac{m_Z^{[3]}}{m_W^{[3]}}
 =x-11\Delta_4+2s_{120}-8\zeta_{270}.
\]
By contrast, the nominal coordinate \((97,9)\) of the scalar sector belongs
to the documented angular chart \(\mathbb Z^2\): this expression does not
attribute to it a \(\mathbb Z^5\) signature or an order-\(3\)
refinement. Its realization as a scalar excitation requires the scalar
projector and the corresponding gauge and fermionic couplings.

For an unstable species, the physical result is a pole of the channel
operator, not merely the evaluation of a positive character.
Subsection~\ref{vi:subsec:resonancias-familias} constructs this reading and
separates pole mass and width. In the preserved documentary register,
the preceding numbers are signature-conditioned evaluations;
the individuation of route, projector, and pole constitutes the sectorial
realization that must accompany them.

\subsection{Finite grammar of material compositions}
\label{vi:subsec:gramatica-compuestos}

Let \(\mathcal Q_a\), \(a\in\mathbb Z/3\mathbb Z\), be the three
color classes of the atlas, each of cardinality \(12\). The formal dual
\(q^\vee\) preserves the conjugate support and reverses charge
and color orientation. It is not merely the cell \(\rho_{\mathrm c}(q)\) stripped
of these data. Mesonic and baryonic expressions are defined by
\begin{align}
 \mathfrak M&=\{q\otimes(q')^\vee:
 q,q'\in\mathcal Q_a\text{ for some }a\},\notag\\
 \mathfrak B&=\mathcal Q_0\times\mathcal Q_1\times\mathcal Q_2.
 \label{vi:eq:gramatica-meson-barion}
\end{align}
The pairs are ordered and the color positions of the triple are
fixed. Residual neutrality is, respectively,
\(a-a=0\) and \(0+1+2=0\pmod3\).

\begin{proposition}[Internal compositional census]
\label{vi:prop:censo-compuestos}
The grammar contains \(432\) mesonic expressions, \(1728\)
baryonic expressions, and \(372\) compatible charged arrows, divided
into \(320\) leptonic and \(52\) quark arrows. These cardinalities count
internal expressions and arrows, not physical species or independent
mass predictions.
\end{proposition}
\begin{proof}
For each color, an ordered pair of \(12\) cells is chosen; therefore,
\(|\mathfrak M|=3\cdot12^2=432\). The baryon's three positions give
\(|\mathfrak B|=12^3=1728\). Lepton--neutrino transitions of
equal depth share only \(G_2,G_4\); in both directions this
gives \(2(8\cdot4+16\cdot8)=320\).

Up--down transitions additionally require equality of color.
In \(G_1\), the color multiplicities are \((2,1,1)\) and
\((0,1,1)\); in \(G_3\), they are \((4,4,4)\) and \((2,2,2)\).
There are \(2+24=26\) pairs in one direction and \(52\) in both. Adding
\(320+52\) completes the census. All factors come from the atlas
preceding mass evaluation.
\end{proof}

The \(81\) cell identities complete the neutral arrows
of this grammar, but a compositional identity does not thereby become
a neutral boson. The representation and realization channel
are the data distinguishing the two notions.

\subsection{Composition of the register and bonding memory}
\label{vi:subsec:composicion-registro}

A composite contains more information than the list of its
constituents. Let \(L_1,L_2\) be enriched registers, \(g\) the
transport aligning phase, sheet, and orientation, and \(B\) the subtrace
identified at the common boundary. In action and
event coordinates, composition is
\begin{equation}
 (n,e)_1\star_{(g,B)}(n,e)_2
 =\bigl(n_1+n_2-n_B,\ e_1+ge_2-e_B\bigr).
 \label{vi:eq:composicion-frontera}
\end{equation}
Subtraction prevents the same subtrace from being counted twice; transport
preserves orientation at the interface. Torsional predecessors and
corona marks are composed at the same level before
\(L_{\mathrm{tick}}\) is applied again. The reversible event decomposition
of Proposition~\ref{vi:prop:eventos} permits the complete register
to be transported without reducing it to its total.

For an ordered tree \(\mathcal T\) of \(n\) constituents, this gives
\begin{equation}
 (\Gamma_1,\ldots,\Gamma_n;\mathcal T,\partial)
 \xmapsto{\mathfrak C_{12}^{(n)}}L_Y
 \xmapsto{L_{\mathrm{tick}}}\sigma_Y
 \longmapsto (\eta_Y,M_Y).
 \label{vi:eq:cadena-compuesto}
\end{equation}
The boundary \(\partial\) includes the transports and subtraces of each
tree edge. Once these admissible data are fixed, each composition
is an explicit operation. The character is applied to the final signature;
it does not preselect the tree through the destination mass.

\begin{proposition}[Compatibility of parenthesization]
\label{vi:prop:cociclo-compuesto}
In a register chart where
\(L_1\star L_2=L_1+L_2+\mathfrak b(L_1,L_2)\), the two
parenthesizations of an admissible triple coincide if and only if
\begin{align*}
 \mathfrak b(L_1,L_2)+\mathfrak b(L_1\star L_2,L_3)
 =\mathfrak b(L_2,L_3)+\mathfrak b(L_1,L_2\star L_3).
\end{align*}
The equality is required on the domain where both trees represent the
same oriented interface.
\end{proposition}
\begin{proof}
Expanding \((L_1\star L_2)\star L_3\) gives
\(L_1+L_2+L_3+\mathfrak b(L_1,L_2)+\mathfrak b(L_1\star L_2,L_3)\).
Expanding \(L_1\star(L_2\star L_3)\) gives the same sum of three
registers and the two terms on the right-hand side. Their equality is
exactly equivalent to the stated identity.
\end{proof}

The cocycle law ensures independence of an equivalent parenthesization;
it does not identify physically different trees. Two bonding structures
preserving different channels remain distinct realizations
until an intertwiner of registers, projectors, and operators is available.

\subsubsection{A finite witness of the information a sum of signatures would lose}

The historical reduced event extractor is defined by
\[
 F_0(n,e)=\left(n,\sum_m e_m,K_0(e),
 \sum_{a=0}^3\operatorname{sgn}(e_a+e_{a+4}+e_{a+8}),
 \sum_m\tau_m\tau_{m+3}\right),
\]
where \(\tau_m=\operatorname{sgn}(e_m)\) and
\(K_0(e)=\sum_{\tau_m=0}\Sigma_{12}(m)\), with
\(\Sigma_{12}=(+,+,+,-,-,-)^2\). Its third coordinate is a
reading of null events; it is not the strong torsion
\(k_{\Delta_4}=T_z-2C_z\) of
Section~\ref{vi:sec:operador-masa}.

Consider the registers of the routes
\[
 a=(1,3,-1,+1),\quad b=(6,7,-1,+1),\quad c=(1,2,-1,-1),
\]
with
\begin{align*}
 e(a)&=(5,4,4,-4,-4,-6)^2,\\
 e(b)&=(4,5,4,-4,-6,-4)^2,\\
 e(c)&=(1,4,1,-2,-4,-3)^2.
\end{align*}
They have \(F_0(a)=F_0(b)=(8,-2,0,0,-12)\) and
\(F_0(c)=(28,-6,0,-4,-12)\). Superposing the registers gives
\[
 F_0(a\boxplus c)=(36,-8,0,0,-12),\qquad
 F_0(b\boxplus c)=(36,-8,0,-2,-12).
\]
The projected inputs are equal and the composite outputs differ.
Consequently, no operation on the two five-tuples of this
projection recovers all superpositions. The example proves a
precise obstruction for \(F_0\); it does not attribute the same result to every
possible extension of a signature. The current construction avoids this
loss through \eqref{vi:eq:cadena-compuesto}: it preserves the events
and predecessors, composes them, and only then evaluates the complete extractor.

\subsection{Mesons: dual pair, singlet, and differentiation of excitations}
\label{vi:subsec:mesones-familias}

The expression \(q\otimes(q')^\vee\) establishes the ordered pair and its
residual neutrality. The color singlet belongs to the realization
\(V\otimes\overline V\), \(V=\mathbb C^3\). If
\(\Omega=\sum_{a=1}^3|a\bar a\rangle\), the projector is
\begin{equation}
 P_M^{(1)}=\frac13|\Omega\rangle\langle\Omega|.
 \label{vi:eq:singlete-meson}
\end{equation}

\begin{proposition}[Mesonic color projection]
The map \(P_M^{(1)}\) is a rank-one orthogonal projector, and
its image is invariant under \(U\otimes\overline U\), \(U\in SU(3)\).
\end{proposition}
\begin{proof}
We have \(\langle\Omega,\Omega\rangle=3\), hence
\((P_M^{(1)})^2=P_M^{(1)}\). Moreover,
\[
 (U\otimes\overline U)\Omega
 =\sum_{a,b,c}U_{ba}\overline U_{ca}|b\bar c\rangle
 =\sum_b|b\bar b\rangle=\Omega.
\]
The image thus preserves the color contraction without privileging any color.
\end{proof}

The mesonic descriptor combines flavors, spin and parity, isospin and charge,
radial or orbital excitation, bonding tree, and boundary. The label
\(J^{PC}\) is used when \(C\) is defined for the channel; charge
conjugation is not automatically assigned to a charged state. Two
states with the same flavors may have different radial,
orbital, or channel projectors and, therefore, different spectra. The
orbital construction of Section~\ref{vi:sec:kepler} provides a
possible reading when its dynamical hypotheses hold; orbital
numbers do not follow from the meson's name alone.

The composite-evaluation register preserves
\[
\begin{array}{c|r r r r r}
 &n_A&s_C&k_{\Delta_4}&b_{120}&b_\zeta\\\hline
 \pi^0&44&-4&-25&1&-2\\
 \pi^\pm&44&1&-2&2&-1\\
 K^\pm&54&-1&17&-2&2\\
 K^0&54&0&32&3&-2
\end{array}
\]
For these signatures, the same reading \eqref{vi:eq:etapas-e3-sectoriales}
explains the differences within the pairs:
\begin{align}
 \log\frac{m_{\pi^\pm}^{[3]}}{m_{\pi^0}^{[3]}}
 &=\frac56y+23\Delta_4+s_{120}+\zeta_{270},\notag\\
 \log\frac{m_{K^0}^{[3]}}{m_{K^\pm}^{[3]}}
 &=\frac16y+15\Delta_4+5s_{120}-4\zeta_{270}.
 \label{vi:eq:desdoblamientos-mesonicos}
\end{align}
The action terms \(44x\) and \(54x\) cancel within
each pair. The remaining differences preserve which register coordinate
produces each contribution. They are exact consequences of the
declared signatures, not a substitute proof of their generation by the
bonding tree.

The resulting evaluations are
\begin{align*}
 m_{\pi^0}^{[3]}&=134.976724327872125739\ldots\ \mathrm{MeV}/c^2,\\
 m_{\pi^\pm}^{[3]}&=139.570485171572191374\ldots\ \mathrm{MeV}/c^2,\\
 m_{K^\pm}^{[3]}&=493.677085649530612475\ldots\ \mathrm{MeV}/c^2,\\
 m_{K^0}^{[3]}&=497.611186497750457358\ldots\ \mathrm{MeV}/c^2.
\end{align*}
Each numerical value retains the
status of an exact evaluation conditioned on the recovered composite
signature. The singlet, grammar, and register extraction are
constructed; complete nominal realization requires the binding
tree to regenerate that signature before comparison begins.

\subsection{Baryons: color triple, exchange, and the neutron--proton difference}
\label{vi:subsec:bariones-familias}

The neutrality of \(q_0\otimes q_1\otimes q_2\) preserves the three
color positions. In the tensor realization, introduce the normalized
vector
\[
 |\varepsilon\rangle=\frac1{\sqrt6}
 \sum_{a,b,c=1}^3\varepsilon_{abc}|abc\rangle,
 \qquad P_B^{(1)}=|\varepsilon\rangle\langle\varepsilon|.
\]

\begin{proposition}[Baryonic singlet and exchange rule]
\label{vi:prop:singlete-barion}
The operator \(P_B^{(1)}\) is a rank-one orthogonal projector,
invariant under \(U^{\otimes3}\), \(U\in SU(3)\). A permutation
of two positions reverses the sign of its generating vector and preserves
the projector.
\end{proposition}
\begin{proof}
There are \(6\) nonzero orthogonal terms, each of modulus
\(1/\sqrt6\); hence the vector has norm one. The determinant
identity gives
\(U^{\otimes3}|\varepsilon\rangle=
\det(U)|\varepsilon\rangle=|\varepsilon\rangle\).
A transposition reverses the sign of \(\varepsilon_{abc}\). In the
outer product of the vector with its adjoint, the two signs cancel.
\end{proof}

Color antisymmetry is one part of the complete state. The physical
projector combines color, spin and parity, isospin, charge, and the exchange
symmetry of orbital and flavor components. Pauli's principle
acts on the complete fermionic product; applying it to a bare
list of colors is insufficient. The triple tree records which pair was coupled
first and how the third constituent crossed the boundary.
Proposition~\ref{vi:prop:cociclo-compuesto} controls equivalent
parenthesizations.

The recovered proton and neutron signatures are
\[
 \sigma_p=(59,0,9,1,0),\qquad
 \sigma_n=(59,1,-34,0,1).
\]
They share the action coordinate \(59\), but differ in sheet,
torsion, and returns. Their exponent difference is
\begin{equation}
 \log\frac{m_n^{[3]}}{m_p^{[3]}}
 =\frac16y-43\Delta_4-s_{120}+\zeta_{270}.
 \label{vi:eq:diferencia-neutron-proton}
\end{equation}
The equality is proved by subtracting the five coordinates. In this signature,
\(-34\) is a torsional integer of the neutron; it is not identified with the
decimal exponent \(-34\) of the absolute action scale. Equality
of two integers in different maps does not eliminate their type or provenance.

The preserved evaluations are
\begin{align*}
 m_p^{[3]}&=938.271751546984898298\ldots\ \mathrm{MeV}/c^2,\\
 m_n^{[3]}&=939.565810472507023312\ldots\ \mathrm{MeV}/c^2.
\end{align*}
Formula \eqref{vi:eq:diferencia-neutron-proton} locates the effect of
each component of the nominal signatures. The physical derivation of the
difference requires isospin, vacancy orientation, the bonding
boundary, and the electromagnetic contribution to be transported into a single chart.
Character evaluation preserves its exactness under the fixed signature;
that exactness does not replace the prospective construction of the bond.

\subsection{Higher-order compositions and channel equivalence}
\label{vi:subsec:compuestos-superiores}

The map \(\mathfrak C_{12}^{(n)}\) admits an arbitrary finite
number of constituents. A tetraquark uses \(qq\bar q\bar q\);
a pentaquark, \(qqqq\bar q\); a gluonic bound state uses
adjoint representations projected onto the singlet; a hybrid preserves
fermionic and gauge constituents in the same tree. Nuclei,
atoms, and molecules use the same precedence of register and reading,
with their own couplings and boundaries.

Binding energy belongs to the realization of the composite register.
It is not obtained by adding isolated masses and naming the desired
residual a binding contribution. Even when a cocycle takes additive values on
registers, the sign and correlation extractor is applied afterwards;
the multiplicative character of a signature therefore does not turn the composite
mass into a sum of constituent masses.

If two trees \(\mathcal T_1,\mathcal T_2\) have a unitary
\(V\) intertwining their operators and channel projectors,
\[
 VH_1=H_2V,\qquad VP_1=P_2V,
\]
then \(V(z-H_1)^{-1}=(z-H_2)^{-1}V\) on the common resolvent set.
Its continuation, when it exists and is transported in the same chart,
preserves poles and their multiplicities. This equality gives the
operational criterion for identifying two realizations of a composite; sharing
flavors or an approximate mass is insufficient.

\subsection{Resonances: pole mass, width, and survival memory}
\label{vi:subsec:resonancias-familias}

Let \(P\) be the prepared subspace and \(Q=I-P\) that of the
eliminated channels. In a region where \(z-QHQ\) is invertible, eliminating
the component \(Q\psi\) from \((z-H)\psi=0\) gives
\begin{equation}
 H_{\mathrm{eff}}(z)=PHP+PHQ(z-QHQ)^{-1}QHP.
 \label{vi:eq:feshbach-familias}
\end{equation}

\begin{proposition}[Reduced reading of a bound or resonant state]
\label{vi:prop:feshbach-familias}
If \(P\mathcal H\) is finite-dimensional, the admissible spectral
values of the reduced problem satisfy
\(\det(z-H_{\mathrm{eff}}(z))=0\). A resonance additionally requires
meromorphic continuation to the relevant channel chart and an isolated
zero with admissible physical residue.
\end{proposition}
\begin{proof}
The equation in the complement is
\((z-QHQ)Q\psi=QHP\psi\). Solving it and substituting into the equation
for \(P\psi\) produces \((z-H_{\mathrm{eff}}(z))P\psi=0\).
In finite dimension, a nonzero solution exists exactly when
the determinant vanishes. The algebraic calculation is performed on the resolvent
domain; passage to a resonance uses the stated continuation and
residue, which do not follow merely from a channel being open.
\end{proof}

A stable state is a normalizable spectral line with no open decay
channel; a bound state lies below the relevant threshold,
although it may decay through another channel; a resonance is a
pole of the continuation. If
\[
 z_X=M_Xc_{\mathrm{int}}^2-\frac{i}{2}\Gamma_X,
 \qquad\Gamma_X>0,
\]
the real and imaginary parts are two coordinated readers of the same
pole. Width is not a sixth coordinate of \(\sigma\).

The discrete reading uses a survival eigenvalue
\(\lambda_X=re^{-i\theta}\), \(0<r<1\), and a time step
\(t_0>0\). Once the branch \(\theta\in I_\theta\) of the quasienergy
chart is fixed,
\begin{equation}
 z_X=\frac{\hbar_{\mathrm{int}}}{t_0}
 (\theta+i\log r),\quad
 M_X=\frac{\hbar_{\mathrm{int}}\theta}{t_0c_{\mathrm{int}}^2},\quad
 \Gamma_X=-\frac{2\hbar_{\mathrm{int}}}{t_0}\log r.
 \label{vi:eq:anchura-discreta-familias}
\end{equation}
Indeed, substitution into
\(\lambda_X=\exp(-iz_Xt_0/\hbar_{\mathrm{int}})\) recovers
\(re^{-i\theta}\). Since \(\log r<0\), the pole is retarded and the
width positive. The survival probability after \(N\) steps is
\(r^{2N}=\exp(-\Gamma_XNt_0/\hbar_{\mathrm{int}})\), whence
\(\tau_X=\hbar_{\mathrm{int}}/\Gamma_X\).

If this reading is obtained from a compression \(SUS\), asymptotic
interpretation through a single mode requires a simple dominant
eigenvalue, nonzero initial overlap, and separation in modulus from the rest
of the spectrum. Normalization on probabilities rather than
amplitudes changes the factor \(2\). The temporal branch and that convention
therefore accompany every expressed width.

Equality of mass and width for conjugate states also requires the
channels to be transported:
\[
 U_C H_{\mathrm{eff}}^X(z)U_C^{-1}=H_{\mathrm{eff}}^{\bar X}(z).
\]
The same poles are then transported with the same causal
prescription. An antiunitary representation additionally requires that
prescription to be transported: bare conjugation of \(z\) interchanges
retarded and advanced poles.

\subsection{Three topological codomains and the virtual sector}
\label{vi:subsec:topologicos-familias}

Topological classes preserve objects and fusion rules, not a
universal scalar mass. The primary datum is
\((N_{ab}^{c},F^{abc}_d,R^{ab}_c)\), with its coherence equations.
An energy gap subsequently requires a Hamiltonian realization and
a spectral projector. Assigning this output type makes it possible
to include these sectors without turning them into elementary-mass rows.

\subsubsection{Fibonacci sector}

The objects \(\mathbf1,\tau\) satisfy
\(\tau\otimes\tau=\mathbf1\oplus\tau\). The positive dimension
\(d_\tau\) preserves fusion multiplication:
\(d_\tau^2=1+d_\tau\). The corpus realization identifies this
dimension with the internal coordinate \(\varphi\) already generated; the
fusion equation is its recognition in this codomain, not a
substitute for its APP--TRIT--TPK genealogy. In the fixed convention,
\[
 F_\tau^{\tau\tau\tau}=\begin{pmatrix}
 \varphi^{-1}&\varphi^{-1/2}\\
 \varphi^{-1/2}&-\varphi^{-1}
 \end{pmatrix},\qquad
 R_{\mathbf1}^{\tau\tau}=e^{-4\pi i/5},\quad
 R_\tau^{\tau\tau}=e^{3\pi i/5}.
\]
The identity \(\varphi^{-2}+\varphi^{-1}=1\) directly proves
that \(F^*F=I\); off-diagonal terms cancel. The phases
of \(R\) preserve the braiding orientation.

\subsubsection{Ising sector and Majorana realization}

The rules are
\[
 \psi\otimes\psi=\mathbf1,\quad
 \psi\otimes\sigma=\sigma,\quad
 \sigma\otimes\sigma=\mathbf1\oplus\psi,
\]
with
\[
 F_\sigma^{\sigma\sigma\sigma}=\frac1{\sqrt2}
 \begin{pmatrix}1&1\\1&-1\end{pmatrix},\qquad
 R_{\mathbf1}^{\sigma\sigma}=e^{-\pi i/8},\quad
 R_\psi^{\sigma\sigma}=e^{3\pi i/8}.
\]
Matrix multiplication gives \(F^2=I\). The sector preserves its
three fusion types and its representation; its Majorana designation
does not automatically turn a neutrino from the preceding sector into a
self-conjugate mode. Such a statement uses the specific antiunitary
conditions of the neutral realization.

\subsubsection{Abelian sector with six classes}

For \(a,b\in\mathbb Z/6\mathbb Z\),
\[
 a\otimes b=a+b\pmod6,\qquad
 R_{ab}=\zeta_6^{ab},\qquad\zeta_6=e^{2\pi i/6}.
\]
The identity \(R_{a+b,c}=R_{a,c}R_{b,c}\) and its analogue in the
second argument express the bilinear character of braiding.
The doubled construction
\(\mathcal C\boxtimes\mathcal C^{\mathrm{rev}}\) preserves both
orientations. Its classes must not be counted as six new elementary
masses.

\subsubsection{Finite persistence and virtuality}

The virtual sector preserves finite extensions, their channel, and their
persistence horizon. A finite section with a terminal obstruction does not occupy
the same codomain as an asymptotically persistent section. Virtuality
is read in an amplitude or resolvent and is not, by
definition, equivalent to a universal imaginary mass. Nor does the absence of
certification of a global extension in a record suffice to assert that
an obstruction exists: the horizon and its proof are different data.
The classification precisely preserves this finite information without
inventing a stable particle for every internal calculation line.

\subsection{Classification of twenty-three particle sectors and scope of their realizations}
\label{vi:subsec:23-clases}

The public inventory comprises \(17\) elementary classes, \(2\) composite
collections, \(3\) topological codomains, and \(1\) virtual sector:
\(17+2+3+1=23\). The three neutrino rows name flavors, the
six quark rows preserve their color triplicity, and the gluon row
preserves the adjoint representation. None of these counting
conventions modifies the atlas \(81/56/13/324\).

The following list retains the output type and the scope of the
documented realizations. The expressions “declared signature” and
“fiber operator” refer to the explicit constructions in this
section, not to different degrees of decimal precision.
The provenance register distinguishes the electronic scalar closure and
the two null charts from physical scalarization still pending in
noncentral collections. For neutrinos, it retains as pending the
complete realization of mixing and scale; for quarks, flavor selection
and its scale chart; for nominal bosons, the joining of route
and pole; for named composites, the scalar realization of their
bond. These documentary statuses do not alter the operators,
grammars, or conditioned evaluations presented above.

\begin{table}[!htbp]
\centering\fontsize{9.2}{10.5}\selectfont
\setlength{\tabcolsep}{4pt}\renewcommand{\arraystretch}{1.04}
\begin{tabular}{@{}p{.11\linewidth}p{.20\linewidth}p{.62\linewidth}@{}}
 \toprule
 Class&Representatives&Preserved result and realization operation\\
 \midrule
 
 1&\(e^-,e^+\)&Rank-one scalar beta closure and conjugate representation.\\
 2&\(\mu^-,\mu^+\)&Exact fiber operators; nominal evaluation and order-\(3\) refinement with a declared signature; distinct physical spectral selection.\\
 3&\(\tau^-,\tau^+\)&Exact fiber operators; nominal evaluation and order-\(3\) refinement with a declared signature; distinct physical spectral selection.\\
 4&\(\nu_e,\bar\nu_e\)&Neutral fiber operator; realization of mixing, mass frame, and physical scale through their complete maps.\\
 5&\(\nu_\mu,\bar\nu_\mu\)&Same neutral domain, with a different flavor label; not an individual atlas depth.\\
 6&\(\nu_\tau,\bar\nu_\tau\)&Same neutral domain; preservation of the Dirac--Majorana alternative under its conditions.\\
 7&\(u,\bar u\)&Color-fiber operator and nominal angular reading; scalarization, scheme, and scale preserved.\\
 8&\(d,\bar d\)&Color-fiber operator and nominal angular reading; scalarization, scheme, and scale preserved.\\
 9&\(c,\bar c\)&Color-fiber operator and nominal angular reading; scalarization, scheme, and scale preserved.\\
 10&\(s,\bar s\)&Color-fiber operator and nominal angular reading; scalarization, scheme, and scale preserved.\\
 11&\(t,\bar t\)&Color-fiber operator and nominal angular reading; explicit mass convention and comparison channel.\\
 12&\(b,\bar b\)&Color-fiber operator and nominal angular reading; scalarization, scheme, and scale preserved.\\
 13&Photon&Exactly zero-mass section; electromagnetic representation.\\
 14&Gluons&Exactly zero-mass section; adjoint color representation.\\
 15&\(W^+,W^-\)&Complete nominal signature and order-\(3\) evaluation; joining to route, projector, and pole as a distinct physical realization.\\
 16&\(Z^0\)&Complete nominal signature and order-\(3\) evaluation; joining to route, projector, and pole as a distinct physical realization.\\
 17&\(H\)&Scalar descriptor and nominal angular evaluation; projector, couplings, and pole of the scalar realization.\\
 18&Mesons&Exact grammar of the dual pair and bond; recovered composite signatures; distinct nominal scalarization and pole channels.\\
 19&Baryons&Exact grammar of the color triple and bond; recovered proton and neutron signatures; distinct nominal binding realization.\\
 20&Fibonacci&Fusion and braiding codomain; a gap requires its Hamiltonian realization.\\
 21&Ising/Majorana&Fusion codomain and real structure according to the representation; no assigned universal scalar mass.\\
 22&\(\mathbb Z/6\mathbb Z\)&Six braiding classes; a topological codomain, not six universal masses.\\
 23&Virtual sections&Finite extension and obstruction horizon; no universal stable scalar.\\
\end{tabular}
\end{table}

The exact documentary statuses are expressed in Appendix~B.3:
Table~\ref{tab:vi-estados-clases} preserves the status legend, and
Table~\ref{tab:vi-clases} preserves, for the twenty-three classes, residence,
incidence, representation, and scalar status.

\begin{proposition}[Realization fiber product and observable codomain]
\label{vi:prop:producto-fibrado-especie-ruta}
For \(X\in\mathcal P_{\rm adm}\), the fiber defined in
\eqref{vi:eq:fibra-descriptor} is nonempty if and only if
\[
 d_{\mathcal P}(X)\in\operatorname{im}d_{\mathcal R}.
\]
On this fiber, the joint spectral measure
\eqref{vi:eq:familias-pvm} determines the observable codomain of the readers
\((M_X^{\beta,D},M_X^{\beta,\Omega})\). A preparation \(P_X\) expresses a
scalar line \(m\) exactly when it satisfies the criterion of
Proposition~\ref{vi:prop:familias-linea}; in all other cases it preserves the
multiplet, spectral distribution, or sectorial type produced.
\end{proposition}
\begin{proof}
By the definition of image, the stated condition is equivalent to the existence of
a route \(\gamma\in\mathcal R_{\rm adm}\) with
\(d_{\mathcal R}(\gamma)=d_{\mathcal P}(X)\); this is precisely the
condition for membership in \(F_X\). The second statement follows from the spectral
theorem applied to the two diagonal operators and the characterization
of single-point support proved in
Proposition~\ref{vi:prop:familias-linea}.
\end{proof}

The external documentary snapshot contains \(471\) mass observables and
\(384\) width observables, totaling \(855\). These are comparison rows, not
\(855\) internal predictions. In particular, the \(243\) mesonic mass
rows and \(166\) baryonic rows do not coincide with the grammar's cardinalities
\(432\) and \(1728\): they count objects of different types.
The external table is opened after identity, fiber, register,
operator, projector, and chart have been fixed.

Reconstruction of a positive ratio through a tower receiving that
ratio demonstrates the representational capacity of the evaluation system.
The independent output of a species, by contrast, requires the tower's
coefficients to arise from its route and boundary. This
separation preserves both representational completeness and effective
sectorial results. The whole is organized by a
common rule: construct the structure and its memory, apply the reader
of the corresponding type, and then perform physical recognition.
